\section{Concluding Remarks}

This work identified vulnerabilities in the \textit{syncfs} system call that expose timing-based leakage vectors, compromising logical isolation. By exploiting these weaknesses, we demonstrated high-bandwidth covert channels. Additionally, we developed three side-channel attacks for fingerprinting websites, videos, and applications. Furthermore, we explored cross-container attacks that breach isolation in containerized environments. Our findings highlight the need for improved defenses, such as redesigning \textit{syncfs} and enhancing I/O buffer isolation. These measures are crucial to mitigating the identified risks while maintaining system functionality.
Our attacks demonstrate that \textit{syncfs}-based vulnerabilities pose a significant threat to system security.

\section*{Acknowledgements}
This research was sponsored in part by the OUSD(R\&E)/RT\&L and was accomplished under Cooperative Agreement Number W911NF-20-2-0267. The views and conclusions contained in this document are those of the authors and should not be interpreted as representing the official policies, either expressed or implied, of the ARL and OUSD(R\&E)/RT\&L or the U.S. Government. The U.S. Government is authorized to reproduce and distribute reprints for Government purposes notwithstanding any copyright notation herein.  We would like to thank the paper reviewer and, especially, the shepherd whose comments and requests substantially improved the paper.

% In this paper, we present covert and side-channel attacks exploiting a novel attack vector based on \textit{syncfs}. By measuring the delay patterns of \textit{syncfs}, we can identify different I/O system call footprints and write sizes. Leveraging these insights, we design a \textit{syncfs} covert channel attack and evaluate its effectiveness on both Linux native file systems and the Windows file system. Additionally, we develop three end-to-end side-channel attacks targeting Linux and Android environments.

% Furthermore, we explore cross-container attacks using \textit{syncfs}, including a container detection technique and a cross-container covert channel attack. Our findings expose a significant vulnerability in \textit{syncfs} and Linux I/O buffers, highlighting the need for enhanced security measures.

% In this paper, we demonstrate covert and side channel attacks using a new attack vector, \textit{syncfs}. In reverse engineering, we identify two leakage vectors through \textit{syncfs}: I/O system call footprint and write size. We utilize these observations to construct a \textit{syncfs} covert channel attack and evaluate it on both Linux native file systems and the Windows file system. We also build three end-to-end side channel attacks on Linux and Android. Additionaly, we present cross-container attacks using \textit{syncfs} including a container detection technique and a cross-container covert channel attack. Our work reveals the a significant vulnerability in \textit{syncfs} and Linux I/O buffers.    